\section{Introdução}\label{sec:intro}

Dado um conjunto de $n$ cidades e as distâncias entre cada par, o Problema do Caixeiro
Viajante (PCV) pede o ciclo de menor custo que visita todas as cidades exatamente uma vez.
Mesmo restrito a pontos no plano com distância euclidiana, o problema continua
NP-difícil~\cite{papadimitriou1977}, e os métodos práticos dividem-se em duas famílias:
algoritmos \emph{exatos}, que garantem o ótimo ao custo de tempo exponencial no pior caso,
e algoritmos \emph{aproximados}, polinomiais e com garantia sobre a qualidade da
solução~\cite{applegate2006}.

Este trabalho implementa e compara três algoritmos: um \emph{Branch-and-Bound} (B\&B)
exato e os aproximados \emph{Twice-around-the-tree} (fator 2) e
Christofides~\cite{christofides1976} (fator 1{,}5). A análise teórica de pior caso,
porém, não diz como esses algoritmos se comportam em instâncias reais de tamanho moderado
nem quanto custam as constantes escondidas na notação assintótica. Por isso o foco é
\emph{empírico}: cada sub-rotina é instrumentada em tempo e memória, e os experimentos
respondem a três perguntas.
\begin{enumerate}
  \item \textbf{Complexidade empírica.} As curvas de tempo e memória medidas seguem os
    modelos assintóticos esperados? Gargalos de implementação, como cópias de matrizes ou o
    coletor de lixo, distorcem essas curvas?
  \item \textbf{Impacto do \emph{warm-start}.} O B\&B começa com um limite superior dado
    por uma rota inicial. Entre cinco estratégias (rota trivial, vizinho mais próximo,
    inserção pelo fecho convexo, Twice-around e Christofides), o tempo gasto na heurística
    compensa a redução da árvore de busca? Onde o método exato colapsa?
  \item \textbf{Sensibilidade.} Que características geométricas da instância
    (distribuição espacial, agrupamento, forma do fecho convexo) pioram ou facilitam a
    qualidade das soluções aproximadas?
\end{enumerate}

Os experimentos usam 15 instâncias euclidianas (\texttt{EUC\_2D}) com 12 a 70 vértices:
as três instâncias completas da TSPLIB~\cite{reinelt1991} nessa faixa e 12
sub-instâncias amostradas de seis instâncias maiores com geometrias diferentes, mais uma
instância de estresse com 75 vértices para localizar o ponto de colapso. Em resumo, os
resultados mostram que:
\begin{itemize}
  \item o B\&B com limite de Held-Karp~\cite{held1970,held1971} prova o ótimo de todas
    as 15 instâncias dentro do limite de 30 minutos, mas nenhuma das cinco estratégias
    fecha a instância com 75 vértices;
  \item a inserção pelo fecho convexo reduz o número de nós expandidos em 77\% (média
    geométrica) em relação à rota trivial; as outras estratégias ajudam pouco, porque o
    primeiro ramo da busca já reproduz a rota do vizinho mais próximo;
  \item o emparelhamento perfeito de custo mínimo responde por até 93\% do tempo do
    Christofides, com crescimento empírico próximo de $n^3$;
  \item a distribuição uniforme dos pontos (entropia espacial e índice de
    Clark-Evans~\cite{clark1954}) piora a qualidade dos algoritmos baseados em árvore
    geradora mínima, mas nenhum modelo único generaliza entre os algoritmos;
  \item na instância mais difícil do conjunto, o limite de Held-Karp fica 2,1\% abaixo do
    ótimo mesmo com o subgradiente convergido: o gap é da própria relaxação.
\end{itemize}

O texto está organizado da seguinte forma. A Seção~\ref{sec:impl} descreve as
estruturas de dados, os algoritmos, a instrumentação e a auditoria do código, gerado com
apoio de uma ferramenta de IA. A Seção~\ref{sec:experiments} apresenta a configuração
experimental, os resultados das três análises e o caso adversarial do limite inferior, e
a Seção~\ref{sec:conclusions} resume as conclusões. O Apêndice~\ref{app:extra} traz
tabelas e figuras complementares. Código, dados brutos, figuras e o histórico de \emph{prompts} estão no
repositório do trabalho.
